\section{Introduction}
\label{sec:introduction}

There is an intriguing phenomenon in public procurement. Some participants in bidding contests lose out systematically for long periods and several times. This fact is interesting because there is a fixed cost for companies to participate in public bids; therefore, under competitive conditions, it would be expected that these companies would win some opportunities or be excluded from this market.

An alternative hypothesis is that these firms, herein referred to as ``frequent losers,'' are operating not autonomously but as part of a cartel in which it is possible to transfer the gains from the cartel among its members. Because of these characteristics, the presence of frequent losers has the potential to serve as a mechanism to detect a cartel in public bids. Usually, antitrust authorities use two main types of cartel detection mechanisms. The first method gained traction throughout the 1990s in the form of leniency agreements and award-winning disclosure \citep{wils2007leniency}. The second mechanism consists of screening methods that became predominant from the mid-2000s due to the greater availability of data and computational capacity. Through observable data such as prices and quantities, screening methods seek to identify behaviors that are consistent with the cartel hypothesis (i.e., coordination among competitors) and inconsistent with the competition hypothesis among participants in a given market.

Although these models might be insufficient to prove a cartel's existence, they are fundamental to drive investigative efforts and may constitute evidence for judicial authorization of search and seizure operations and interception of communications.

Screening methods available in the literature, such as \citet{green1984noncooperative}, \citet{haltiwanger1991the}, and \citet{abrantesmetz2006a}, have two fundamental limitations. First, these models cannot distinguish tacit collusion from cases of explicit collusion (i.e., cartels). This limitation is problematic for their practical application since tacit collusion, unlike cartels, does not constitute an antitrust offense \citep{harrington2008detecting}.

Second, screening methods solely identify collusive behavior after its occurrence and therefore serve as support only for the repression of cartels but not for their prevention. In public procurement, it would be desirable for screening methods to identify collusive behavior before the bidding process occurs, through observable variables in public notices and participants' registration, to reduce the social costs of cartels.

This paper proposes an alternative screening method that directly addresses these two limitations and does not require additional data beyond those already available to competition authorities and control bodies.

Primarily, the proposed method consists of the use of frequent losers as ``flags'' for the ex-ante screening of cartels in public tenders. To this end, this research proposes a method for identifying frequent losers to differentiate these companies from others that, although they may lose frequently, do not exhibit sufficiently abnormal behavior. The article also presents frequent losers' descriptive characteristics, reinforcing their properties as markers of collusion.

Finally, this paper examines the relationship between the presence of frequent losers in public bids and the price level and other indicators of competition intensity, such as the number of participants and the number of bids. The results indicate that tenders in which frequent losers participate have 7\% higher prices, a 54\% higher number of participants, and a 52\% increase in the number of bids.

These results are consistent with cartel behavior (higher prices) and a strategy to avoid detection by increasing signs of higher competitive intensity. Note that this is precisely the expected effect of frequent losers: more participants and no competitive pressure on the winning bid. External validation using CADE cartel convictions in public procurement confirms that frequent losers serve as \textit{cover bidders}---appearing alongside convicted cartel firms at 2.6 times the baseline rate---rather than being the cartel ringleaders themselves.

This paper is organized as follows. Section~\ref{sec:literature} reviews the related literature on cartel screening methods. Section~\ref{sec:theory} develops the theoretical framework and testable predictions. Section~\ref{sec:data} presents the data and the method for identifying frequent losers. Section~\ref{sec:empirical_strategy} describes the empirical strategy. Section~\ref{sec:results} presents the main results. Section~\ref{sec:robustness} provides robustness checks and additional tests. Section~\ref{sec:conclusion} concludes.
